% Work record, not a completed implication from Sun's conjecture.
\paragraph{The exact good-modulus calculation.}
For $q\geq1$ and $c=(c_1,c_2,c_3)\in\Z^3$, define
\[
 S_q(c)=\sum_{a\bmod q}^{*}\sum_{b\in(\Z/q\Z)^3}
 e\left(\frac{a(b_1^2-3b_2^2-18b_3^2+2)+c\cdot b}{q}\right).
\]
If $\gcd(q,6)=1$, completion of the three squares gives exactly
\begin{equation}\label{A5audit:ternary-gauss}
 S_q(c)=\left(\frac6q\right)G(1;q)^3
 T(2,\overline{72}H(c);q),
 \quad H(c)=-18c_1^2+6c_2^2+c_3^2,
\end{equation}
where $G(1;q)=\sum_{y\bmod q}e(y^2/q)=\epsilon_q\sqrt q$.
Indeed for $d\in\{1,-3,-18\}$,
\[
 \sum_{b\bmod q}e((ad b^2+c_i b)/q)
 =e(-\overline{4ad}c_i^2/q)
  \left(\frac{ad}{q}\right)G(1;q).
\]
The product of the three Jacobi factors is
$\bigl(\frac a q\bigr)\bigl(\frac{54}q\bigr)
=\bigl(\frac a q\bigr)\bigl(\frac6q\bigr)$, and the completed
phase is $2a+\bar a\overline{72}H(c)$.
The symmetry $T(m,n;q)=T(n,m;q)$ follows by $x\mapsto\bar x$.
This proves \eqref{A5audit:ternary-gauss} with all signs fixed.
The program \path{code/B_A5_delta_audit.py} verifies it in the exact
cyclotomic ring for 60 pairs, including prime-power and composite
moduli and both signs of $H(c)$.

This calculation identifies the obstacle to an immediate application
of Steiner's integral-weight conjecture: the Jacobi character in the
inner modulus sum is present.  To use Sun's half-integral-weight
Conjecture~1.4 one must still match the sums, including their local
factors at $2$ and $3$, to precisely specified multiplier systems on
$\Gamma_0(N)$ and verify the stated spectral condition there.
No such complete bad-prime and multiplier identification is claimed
in this record.

\paragraph{Where a sharp region introduces a further uniformity problem.}
In the usual delta-symbol expansion with parameter $X$, the
oscillatory integrals have the form
\[
 I_q(c)=\int_{\R^3}w(x/X)
 h\left(q/X,(F(x)+2)/X^2\right)e(-c\cdot x/q)\,dx,
\]
with $F(s,t,z)=s^2-3t^2-18z^2$.
The sharp physical region is a rational tetrahedral cone cut by a
height plane; its top vertices include
$(4,0,1)$, $(17/4,0,1)$ and $(17/4,1/4,1)$ in normalized coordinates.
Integration by parts along its faces produces inverse-modulus phases
$e(-Xc\cdot v/q)$ at vertices $v$.
For example the upper $t$ boundary is $t=s-4z$, so when
$c=(0,k,0)$ its boundary phase is $e(-k(s-4z)/q)$; the upper
$s$ and $z$ endpoints then give $e(-kX/(4q))$.
This example identifies phases in the direct full-tetrahedron
organization; it does not assert that this boundary term alone is a
lower bound for the actual arithmetic error.  An alternative extension
of the cutoff away from the quadric can suppress vertices not lying
near $F=0$.  Such a reorganization must be analyzed separately; we do
not assert that these particular vertex terms are unavoidable in every
delta-method formula for the same count.

For a dyadic modulus block $Q<q\leq2Q$, partial summation of an
untwisted sum against $e(D/q)$ costs total variation of order
$1+|D|/Q$.  With $Q\asymp X$, $D\asymp X|c|$, and dual
frequencies as large as $|c|\asymp X^{1/2}$ after smoothing at
physical relative width $X^{-1/2}$, this is a polynomial loss.
Sun's conjecture, which bounds the untwisted partial sums, does
not itself state the needed bound after such a phase is inserted.
No free suppression of that variation is valid.

The latter warning has an exact abstract model.  For an integer
$Q\geq16$ and $D=Q^{3/2}$, put
\[
 a_q=e(-D/q)-e(-D/(q-1)),\qquad Q<q\leq2Q.
\]
Every interval sum of $a_q$ has modulus at most $2$, by telescoping,
and $|a_q|\leq2\pi D/(Q(Q+1))\ll Q^{-1/2}$.
Nevertheless
\begin{align*}
 \operatorname{Im}\sum_{Q<q\leq2Q}a_q e(D/q)
 &=\sum_{Q<q\leq2Q}\sin\frac{2\pi D}{q(q-1)}\\
 &\geq4D\sum_{Q<q\leq2Q}\frac1{q(q-1)}
 =\frac{2D}{Q}=2\sqrt Q.
\end{align*}
Here every sine argument lies in $[0,\pi/2]$, and
$\sin t\geq2t/\pi$ on that interval.
Thus even an $O(1)$ untwisted partial-sum bound together with the
individual size $O(Q^{-1/2})$ does not imply an
$O(Q^\varepsilon)$ bound after the relevant inverse-modulus twist.
Arb evaluations for $Q=16,100,10000$ are preserved in
\path{code/B_A5_delta_audit.out}.
These $a_q$ are an abstract example, not Kloosterman coefficients.
The example identifies a gap in the suggested partial-summation
inference; it does not prove that Sun's full arithmetic conjecture
can never imply the desired count by another argument.

\paragraph{Present status.}
An exact follow-up now tests the full amplitude, rather than the phase
alone; see \path{records/B_A5_compensated_weight.tex}.  Multiplication
by a smooth function of $F(y)$ that equals one on every counted point
for $X\geq16$ makes the actual Heath-Brown kernel equal to
$\omega(r)/r$ on $r=q/X\in[2/3,3/4]$.  For $c=(1,0,k)$ the full
normalized Abel weight has size $\asymp k^{-1}$, but its total
variation is $L+O(k^{-1})$, where $L>0$ is explicitly certified.
The summed discrete variation through $H$ is
$LH+O(\log(2H)+H^2/X)$, and boundary smoothing at width $H^{-1}$
retains a linear lower bound on the summed variation for
$k\leq\eta H$.  Thus the amplitude compensates for one phase factor
without giving an epsilon-sized sum of variation norms.  This is
still only an obstruction to the termwise Abel/absolute-frequency
procedure, not a lower bound for an arithmetic error or an
impossibility assertion about all deductions from Sun's conjecture.

The next required step is either a new uniform transformation which
absorbs these phases into the exact half-integral Kloosterman sums,
or a published conjecture which already covers them with its theta
terms removed and its dependence on $D$ explicit.  We have not
found and verified such an implication.  Replacing Sun's statement
by a stronger twisted statement without identifying a primary
source would invent the requested conjectural input.
